\documentclass[11pt,a4paper]{article}
\frenchspacing
\usepackage[margin=22mm]{geometry}
\usepackage{booktabs,tabularx,array}
\usepackage[hidelinks]{hyperref}
\setlength{\parskip}{4pt}
\setlength{\emergencystretch}{2em}

\newcommand{\file}[1]{\texttt{#1}}

\title{Comparing Claude, GPT, and Qwen Treatments\\of Pathfinder Campaigns}
\author{Pathfinder campaign review}
\date{19 September 2026}

\begin{document}
\maketitle

\begin{abstract}
This note compares the completed original Claude campaign with later
GPT-5.6-sol and Qwen3.5-397B full-protocol campaigns.  It combines descriptive
statistics with close reading of matched paper pairs.  Claude behaved most
like an autonomous research programme: it derived, simulated, inspected
sources, and frequently rejected its own proposed connexions.  GPT gave the
strongest compact treatment from abstract-level evidence, especially in
causal framing and experiment design, but sometimes promoted a polished
proposal as a result.  Qwen was effective at detecting structural mismatches
and triaging ideas, but usually stopped before producing a precise estimand,
theorem, or empirical result.  These differences are informative but not a
clean model benchmark because the Claude campaign had more rounds, peer
research, and source access.
\end{abstract}

\section{Scope and comparability}

The comparison uses 14 completed threads in the original Claude directory,
12 in the GPT full-protocol campaign, and 10 valid threads in the Qwen
full-protocol campaign.  Qwen threads Q4P3 and Q4P9 are excluded throughout:
their own artifacts report that peer work terminated before the planned
protocol completed.  Four pairs appear in all three valid sets: Q3P3, Q3P10,
Q4P6, and Q4P10.  These matched pairs support direct qualitative
comparison.  Campaign-wide counts describe each realised run but also reflect
different shortlists.

The treatment was unequal.  The original campaign commonly allowed two or
three rounds, peer research, full-text or code inspection, and verifier-led
continuations.  The GPT and Qwen campaigns used two seats, one round, no peer
search, and predominantly abstract-only inputs.  Selection also differed:
for example, GPT ranked Q4P3 first in its shortlist, while Qwen ranked it last
among its twelve selections.  The statistics below are therefore descriptive
diagnostics, not causal estimates of model quality.

\section{Summary statistics}

\begin{table}[ht]
\centering
\small
\begin{tabular}{lrrr}
\toprule
 & Claude & GPT & Qwen \\
\midrule
Valid completed threads & 14 & 12 & 10 \\
Final DRAFT & 5 & 6 & 3 \\
Final PAUSE & 8 & 6 & 0 \\
Final ITERATE & 1 & 0 & 7 \\
DRAFT rate & 35.7\% & 50.0\% & 30.0\% \\
Ledger entries, total & 471 & 262 & 49 \\
Ledger entries, median per thread & 33.5 & 21.5 & 4.5 \\
Ledger entries, mean per thread & 33.6 & 21.8 & 4.9 \\
Final note lines, median & 70.0 & 33.0 & 55.5 \\
Final note lines, mean & 141.4 & 39.2 & 55.6 \\
Peer research files under \file{ada/} or \file{emmy/} & 59 & 0 & 0 \\
\bottomrule
\end{tabular}
\caption{Campaign-wide descriptive statistics after excluding known faulty-
harness threads.  PAUSE-ON-ITERATE is counted as the final ITERATE verdict.
File and line counts measure realised campaign activity, not intrinsic model
capacity.}
\end{table}

The scale differences are large.  Claude produced about 1.5 times as many
ledger entries per thread as GPT and about 6.9 times as many as Qwen.  GPT's
median final note was half the length of Claude's, despite retaining a much
larger ledger than Qwen.  Claude's very high mean note length relative to its
median reflects several unusually deep investigations.  The 59 research files
are the clearest operational difference: Claude generated simulations,
analysis scripts, numerical outputs, theorem notes, and source inspections,
whereas the one-round GPT and Qwen campaigns generated no corresponding peer
research directories.

Final decisions also differed.  GPT split evenly between DRAFT and PAUSE.
Qwen kept seven of ten valid threads alive as ITERATE and drafted three.  Claude
drafted five, paused eight, and left one at ITERATE.  Thus Qwen was not simply
``more conservative'': it was less decisive, preferring another check to a
terminal rejection.  GPT was most willing to treat an identification-ready
design as draft-worthy.  Claude usually demanded a theorem, measurement, or
other supported pair-level result.

\section{Qualitative model signatures}

\subsection{Claude: sustained research and falsification}

Claude treated each pair most like an autonomous research programme.  It
developed formal results, wrote and ran simulations, inspected released code,
searched prior work, and tracked objections through correction or withdrawal.
Its strongest habit was epistemic separation: association versus intervention,
selection versus treatment, and measurement versus mechanism were kept
distinct.  It was willing to derive a substantial result about one paper and
still pause because the cross-paper bridge remained unsupported.

Its weakness was cost and scope.  Some treatments became long technical
investigations whose strongest contribution no longer belonged to the pair.
The resulting notes were the most precise and auditable, but also the least
accessible and the most expensive to produce.

\subsection{GPT: compact causal design}

GPT was the strongest one-round, abstract-only researcher.  Its characteristic
move was to reject a broad analogy, isolate the identification problem, define
matched treatment arms and estimands, list confounds, and state exactly what
the available evidence could not establish.  The writing was concise,
structured, and publication-oriented.  Relative to Qwen, GPT more often
specified factorial designs, primary and secondary outcomes, non-inferiority
criteria, and construct-validity checks.

Its main weakness was an inconsistent acceptance threshold.  Several threads
were promoted to DRAFT because they contained a rigorous mechanism-discrimination
proposal or experimental design, even though the note itself acknowledged
that no effect had been measured.  GPT therefore sometimes treated a strong
research plan as a research result.

\subsection{Qwen: conceptual triage}

Qwen was effective at identifying broad structural mismatches: false consensus
versus market heterogeneity, cooperative credit assignment versus strategic
concealment, committed-principal mechanisms versus bilateral negotiation, and
individual pathology versus market-level outcome.  This made it useful as an
early critic and explains the high ITERATE rate.

It usually stopped one level too early.  Its next step was often to inspect
traces, formalise a model, or search for sandbagging without specifying the
decisive estimand, counterfactual, theorem assumptions, or controlled contrast.
It occasionally upgraded structural resemblance into direct applicability
before verifying that the source theorem survived the reframing.  Several TeX
files also retained Markdown fences.  The two known incomplete peer runs are
excluded rather than interpreted as model behaviour.

\section{Matched cases}

\subsection{Q3P3: adversarial review and double auctions}

Claude reframed the pair around institutions and stopping rules, implemented
constrained zero-intelligence auction simulations, decomposed efficiency loss
into omission and commission, and inspected source code for constraint
differences.  It pursued a separate continuation and nevertheless paused when
the matched code-review ablation could not be executed from the available
materials.  See \file{threads/Q3P3/Q3P3.tex} and
\file{experiments/2026-09-16-Q3P3-continuation/threads/Q3P3/}.

GPT diagnosed the proposed intervention as underidentified: adding adversarial
review changes deliberation, compute, tokens, delay, and disagreement at once.
It proposed a matched naive-critic versus evidence-grounded-critic contrast,
with allocative efficiency as the primary estimand, and paused because this
was only a research design.  See
\file{experiments/2026-09-18-gpt-5-6-sol-full-protocol/threads/Q3P3/Q3P3.tex}.

Qwen identified the clean conceptual objection: the auction paper reports
heterogeneity, whereas adversarial review addresses false consensus.  It
suggested examining chain-of-thought traces for premature commitment, but did
not develop an identification strategy or perform an empirical test.  See
\file{experiments/2026-09-18-qwen3-5-397b-full-protocol/threads/Q3P3/Q3P3.tex}.

\subsection{Q4P6: strategic credit assignment}

Claude derived a population identification theorem.  With a shared latent
comparison and conditionally independent symmetric observer noise,
cross-observer moments identify observer reliability; calibrated first moments
then identify comparison probabilities, and a connected Bradley--Terry graph
identifies contribution scores up to translation.  It explicitly separated
this result from unproved finite-sample, collusion, bootstrap, and churn claims.
See \file{threads/Q4P6/Q4P6.tex}, especially lines 32--59.

GPT produced a useful but shallower separation result: report-independent
potential shaping cannot repair profitable sandbagging, while a transfer-based
repair requires full-policy implementability constraints and identifiable
signals.  It also made the key information-theoretic point that no mechanism
can distinguish a sandbagger from a genuinely low-capability agent if they
induce the same comparison distribution.  See
\file{experiments/2026-09-18-gpt-5-6-sol-full-protocol/threads/Q4P6/Q4P6.tex}.

Qwen recognised the cooperative-versus-strategic mismatch and proposed looking
for sandbagging in the experimental tasks.  It did not derive a mechanism,
identification condition, or theorem, and returned ITERATE.  See
\file{experiments/2026-09-18-qwen3-5-397b-full-protocol/threads/Q4P6/Q4P6.tex}.

\section{Interpretation}

The models appear to optimise for different stopping questions:
\begin{itemize}
  \item Claude asks whether a demonstrated new pair-level result exists.
  \item GPT asks whether a defensible and identifiable research contribution
        has been specified.
  \item Qwen asks whether a plausible connexion remains worth one more check.
\end{itemize}

This explains more of the verdict pattern than a simple optimistic-versus-
pessimistic scale.  Claude can invest heavily and still reject.  GPT can draft
a proposal while carefully denying that it is an established empirical effect.
Qwen can identify a serious mismatch yet preserve the thread as ITERATE.

\section{Operational conclusion}

For Pathfinder, the models are potentially complementary.  Qwen is useful for
cheap initial mismatch screening.  GPT is the best compact option for narrowing
a connexion into an identifiable theorem or experiment.  Claude is strongest
for expensive follow-through involving formal derivation, source auditing, or
empirical work.  A model-independent verifier should enforce one acceptance
criterion across all three: a polished proposal, a negative conceptual
observation, and a demonstrated pair-level result should not all receive the
same DRAFT status.

The principal experimental follow-up is a controlled rerun with the same
shortlist, source package, number of rounds, peer-search permissions, tool
budget, and verifier.  Until then, the present results support qualitative
workflow choices, but not a causal ranking of intrinsic model quality.

\end{document}
